\section{Literature Review}
\label{sec:literature-review}

\subsection{Researcher Degrees of Freedom and Specification Surfaces}

Empirical work often foregrounds a preferred specification, with robustness
checks reported as supporting evidence. The inferential difficulty is that the
preferred specification is often itself the product of choices made after
examining the data. Choices of outcomes, samples, transformations, exclusions,
controls, stopping rules, and reporting rules can then make nominal tests
behave differently from their advertised level
\citep{simmonsFalsePositivePsychologyUndisclosed2011,johnMeasuringPrevalenceQuestionable2012,wicherts2016degrees}.
The literatures on HARKing and the garden of forking paths emphasize that this
need not look like an explicit search: a single reported analysis can be
data-contingent because the route to it was shaped by the realized data
\citep{kerr1998harking,gelmanLoken2013,gelman2014crisis}.

We call the object over which such choices range a \emph{specification
surface}. Its domain is the set of admissible specifications, preprocessing
choices, or representations available to the analyst. Its value at each point
is the reported quantity of interest, such as a coefficient, standard error, or
test statistic. This terminology separates the object from any particular
display: a specification curve plots part of a surface, while a multiverse
analysis enumerates points on it.

Several literatures study related manifestations of this problem. Publication
bias, file-drawer effects, bunching or discontinuities in reported test
statistics, and p-curve diagnostics study selection across studies or reported results
\citep{rosenthal1979file,franco2014publication,brodeur2016star,head2015extent,brunsIoannidis2016pcurve,brodeur2020methods,vivalt2019specification,jerke2025publication};
recent econometric work develops corresponding diagnostics and corrections
\citep{andrews2019identification,elliott2022detecting}. Model-uncertainty and
specification-search work shows that conclusions can vary sharply across
reasonable models
\citep{leamer1983con,salaimartin1997twoMillion,young2009modelUncertainty,youngModelUncertaintyRobustness2017}.
Recent computational work on coefficient stability is closer still:
\citet{knaeble2024branch} search for maximum and minimum adjusted slope
coefficients over a large discrete model space using branch-and-bound methods.
That line treats the finite optimization problem directly; our goal is to
characterize the geometry and limiting size of the envelope.
Specification-curve, multiverse, and many-analysts designs make this variation
visible by replacing a single preferred model with a structured collection of
defensible alternatives
\citep{steegenIncreasingTransparencyMultiverse2016,simonsohnSpecificationCurveDescriptive2019,silberzahn2018many,breznau2022hidden,rohrer2026multiverse}.
Our question is more local. Conditional on one dataset and one specification
surface, how large can the reported coefficient or $t$-statistic become?

Pre-registration and pre-analysis plans address the same concern by restricting
the specification surface before the data are analyzed
\citep{nosek2018preregistration,olken2015promises,banerjee2020praise,kasy2021forking}.
Empirically, such plans appear most informative when they are sufficiently
specific \citep{brodeurPreregistrationPreanalysisPlans2024}. Theoretically,
they are connected to broader models of scientific communication, transparency,
p-hacking incentives, and implementable statistical decisions under misaligned
researcher and social objectives
\citep{andrews2021model,libgober2022false,mccloskey2022incentive,spiess2018optimal,kasy2022optimal}.
These approaches constrain the surface. They do not characterize the geometry
of the remaining surface, or the extreme values that can be obtained by
searching it.

\subsection{Experimental Specification Surfaces}

Randomized experiments give an especially clean case. Randomization fixes the
treatment direction, and covariate adjustment is not needed for identification;
it is really just a precision device \citep{lin2013agnostic}. More generally, the
assignment mechanism, and hence the propensity score, is known by design. In
high-dimensional semiparametric problems, this information can change what is
statistically possible. This is the lesson of the curse-of-dimensionality
appropriate asymptotics of \citet{robinsRitov1997coda}, and it is emphasized
for randomized controlled trials by
\citet{aronowRobinsSaarinenSavjeSekhon2021}.

In our setting, the implication is geometric. In an experiment, moving across
the specification surface does not change the population target of the
treatment coefficient. It can still change the reported coefficient or
studentized statistic, but the movement comes from chance alignments between
the randomized treatment vector and candidate adjustment directions. This is
the regime in which a randomization-based or reduced-form approximation to the
searched statistic is most natural.

\subsection{Observationally Confounded Specification Surfaces}

Observed confounding changes the surface. In observational work, controls are
often part of the identifying argument. Adding, removing, or transforming them
can make conclusions model-dependent and can change the estimand, the residual
treatment variation, and the apparent strength of evidence
\citep{kingDetectingModelDependence2007,lenz2021achieving,sturman2022uncontrolled,yu2024misstatements,cinelliForneyPearl2022crash}.
Thus an observationally confounded specification surface is not merely an
experimental surface with a different variance.

The closest conceptual precedents are omitted-variable-bias and sensitivity
analysis. These formulas decompose the effect of an omitted covariate into its
association with treatment and its association with the outcome
\citep{oster2019unobservable,cinelliHazlett2020sensitivity,chernozhukov2022long};
related work extends the same style of sensitivity analysis to instrumental
variables \citep{cinelliHazlett2025iv}.
\citet{tsao2025minimum}, for example, ask how strong an omitted covariate
would need to be to overturn an insignificant result. We use the same two
associations for a different purpose: the covariates are observed, and the
analyst moves across the specification surface by deciding which directions to
adjust for or how to represent them. Recent design-phase work similarly treats
covariate balance and subsample construction as ways to reduce exposure to
misspecification \citep{choi2025designPhase}. Here the design choice itself is
part of the specification surface, and the question is how much movement remains
available after the admissible surface has been fixed.

\subsection{Post-Selection and High-Dimensional Inference}

Post-selection inference constructs valid tests or confidence sets after a
specified selection procedure, using tools such as data splitting and selective
conditioning
\citep{cox1975dataSplitting,berk2013valid,fithian2014optimal,lee2016exact,kuchibhotla2018assumptionLean,kuchibhotlaKolassaKuffner2022postSelection,tianTaylor2018selective}.
\citet{berk2013valid} is particularly close in spirit because their
post-selection-inference constant controls a maximum over all submodels,
uniformly over selection rules. The object differs from ours: that framework
targets finite-sample validity for studentized statistics after
coordinate-submodel selection, while we study the
coefficient envelope over a specification surface and its drift, variance, and
outcome-energy decomposition.

High-dimensional regression and proportional-asymptotic random-matrix theory
develop related methods when models or nuisance components are estimated from
many covariates
\citep{el2013robust,sur2019likelihood,hastie2022surprises}. Those results give
the calculus for regimes with $p/n$ non-negligible, but their standard targets
are fixed procedures, likelihood ratios, or prediction risk rather than a
search-maximized reported coefficient. In the low-dimensional sparse case, our
rates also draw on the restricted-isometry and restricted-eigenvalue geometry
used in compressed sensing and sparse regression
\citep{candesTao2005decoding,baraniuk2008simple,bickelRitovTsybakov2009,vandegeerBuhlmann2009conditions,raskuttiWainwrightYu2010restricted}.
Our use is reversed: the same geometry that enables sparse recovery also
controls how much an adversarial specification search can extract.

Unlike these works, we do not condition on a single observed
selection event, and we do not aim to recover a sparse true model. We study the
envelope of possible reported results over a specification surface when the
search rule may be informal, hidden, or adversarial.

We measure researcher discretion through geometry rather than
through a count of possible regressions. Experimental search,
observationally confounded search, and general preprocessing define different
specification surfaces, with different extreme-value behavior. The relevant
degrees of freedom are determined by the covariance geometry of treatment,
outcome, and the admissible representations available to the analyst.
